\section{Geometry, Projection \& Classical Vision}\label{sec:geom}
\lab{epipolar geometry \textperiodcentered\ pinhole optics \textperiodcentered\ feature detectors \textperiodcentered\ filtering \textperiodcentered\ optical flow}

\noindent\begin{minipage}[t]{0.44\columnwidth}\vspace{0pt}\centering
\resizebox{\linewidth}{!}{%
\begin{tikzpicture}[x=1mm,y=1mm,every node/.style={font=\tiny}]
\draw[draw=acc,thick] (0,0) -- (12,8) -- (12,-8) -- cycle;
\node[acc,left] at (0,0) {$O$ (pinhole)};
\draw[draw=mute,dashed] (0,0) -- (32,0);
\draw[draw=acc2,fill=soft2] (18,-10) rectangle (20,10);
\node[acc2,above] at (19,10) {image plane};
\draw[arr,acc3] (0,0) -- (28,7);
\fill[acc3] (28,7) circle (0.6mm) node[above right] {$X \in \mathbb{R}^3$};
\fill[acc] (19,4.75) circle (0.5mm) node[below right] {$x \in \mathbb{R}^2$};
\draw[<->,mute] (0,-12) -- (19,-12) node[midway,below] {$f$ (focal)};
\end{tikzpicture}%
}
\end{minipage}\hfill
\begin{minipage}[t]{0.54\columnwidth}\vspace{1pt}\small
\textbf{World to pixel coordinates.} 3D world point $\mathbf{X}_w = [X, Y, Z, 1]^\top$ maps to 2D homogeneous pixel $\mathbf{x} = [u, v, 1]^\top$ via:
$$\lambda \begin{bmatrix} u \\ v \\ 1 \end{bmatrix} = \mathbf{K} [\mathbf{R} \mid \mathbf{t}] \begin{bmatrix} X \\ Y \\ Z \\ 1 \end{bmatrix}$$
where $\mathbf{K}$ is the $3\times 3$ intrinsic matrix, $\mathbf{R} \in SO(3)$ is rotation, and $\mathbf{t} \in \mathbb{R}^3$ is translation.\\[2pt]
{\color{acc2}$\blacktriangleright$} Intrinsic matrix: $f_x, f_y$ focal lengths, $c_x, c_y$ principal point, $s$ skew.\\
{\color{acc2}$\blacktriangleright$} Depth ambiguity: $\lambda = Z_c$ is lost in projection.
\end{minipage}

\begin{mem}[Epipolar Geometry \& Core Transformations]
Two calibrated cameras observing scene point $X$ have image coordinates $x, x'$.
\begin{itemize}
\item \textbf{Essential Matrix} $\mathbf{E} = [\mathbf{t}]_\times \mathbf{R} \in \mathbb{R}^{3\times 3}$: relates normalized camera coordinates: $x'^\top \mathbf{E} x = 0$. $\mathbf{E}$ has 5 degrees of freedom (two equal singular values, one zero).
\item \textbf{Fundamental Matrix} $\mathbf{F} = \mathbf{K}'^{-\top} \mathbf{E} \mathbf{K}^{-1}$: relates uncalibrated pixel coordinates: $x'^\top \mathbf{F} x = 0$. $\mathbf{F}$ has rank 2, 7 degrees of freedom. Epipolar line $l' = \mathbf{F} x$. Epipoles satisfy $\mathbf{F} e = 0, \mathbf{F}^\top e' = 0$.
\item \textbf{Normalized 8-Point Algorithm}: Hartley's algorithm translates point centroid to $(0,0)$ and scales average distance to $\sqrt{2}$ before constructing matrix $A f = 0$. Enforces $\det(\mathbf{F})=0$ by zeroing smallest singular value in SVD.
\item \textbf{Homography} $\mathbf{H} \in \mathbb{R}^{3\times 3}$: maps coplanar 3D points between two views: $x' \sim \mathbf{H} x$. Solved with 4 point pairs using Direct Linear Transformation (DLT) + RANSAC.
\end{itemize}
\end{mem}

\begin{qa}[Classical Features, Edges \& Keypoints]
\Q{Walk through the 4 stages of Canny edge detection in detail.}{
\textbf{1. Gaussian smoothing}: Convolve image with Gaussian filter ($G_\sigma, \sigma \approx 1.4$) to suppress high-frequency sensor noise.
\textbf{2. Gradient calculation}: Apply Sobel operators $G_x = \begin{bmatrix}-1&0&1\\-2&0&2\\-1&0&1\end{bmatrix}, G_y = G_x^\top$ to compute gradient magnitude $|G| = \sqrt{G_x^2 + G_y^2}$ and direction $\theta = \arctan(G_y / G_x)$. Quantize $\theta$ into 4 canonical sectors: $0^\circ, 45^\circ, 90^\circ, 135^\circ$.
\textbf{3. Non-Maximum Suppression (NMS)}: Thin wide edges down to 1-pixel ridges. At each pixel, compare $|G|$ with its two neighbors along direction $\theta$. If $|G|$ is smaller than either neighbor, suppress it to 0.
\textbf{4. Double thresholding \& Hysteresis}: Define high threshold $T_{\text{high}}$ and low threshold $T_{\text{low}}$ ($T_{\text{high}} \approx 2-3\times T_{\text{low}}$). Pixels with $|G| \ge T_{\text{high}}$ are marked as strong edges; pixels with $T_{\text{low}} \le |G| < T_{\text{high}}$ are weak edges. Retain weak edges \emph{only} if connected to a strong edge via 8-connectivity.}
\Q{How does the Harris Corner Detector mathematically separate corners, edges, and flats?}{
Compute second moment matrix $M = \sum_{(u,v)} w(u,v) \begin{bmatrix} I_x^2 & I_x I_y \\ I_x I_y & I_y^2 \end{bmatrix}$, where $w(u,v)$ is a Gaussian window.
The eigenvalues $\lambda_1, \lambda_2$ of $M$ characterize the local intensity curvature:
\begin{itemize}
\item \textbf{Flat region}: $\lambda_1 \approx 0, \lambda_2 \approx 0$ (no intensity variation).
\item \textbf{Edge}: $\lambda_1 \gg \lambda_2 \approx 0$ (variation along gradient normal, invariant along tangent).
\item \textbf{Corner}: $\lambda_1, \lambda_2$ both large (variation in all directions).
\end{itemize}
Harris avoids explicit eigendecomposition by defining corner response $R$:
$$R = \det(M) - k \cdot (\text{Tr}(M))^2 = \lambda_1 \lambda_2 - k (\lambda_1 + \lambda_2)^2 \quad (k \in [0.04, 0.06])$$
$R > 0$ indicates a corner, $R < 0$ indicates an edge, and $|R| \approx 0$ indicates a flat region.}
\Q{Compare SIFT, SURF, and ORB feature extractors.}{
\textbf{SIFT (Lowe, 2004)}: Scale-space DoG (Difference of Gaussians) extrema localization; 3D quadratic sub-pixel refinement; gradient orientation histogram assignment (rotation invariance); $4\times 4$ spatial grid $\times 8$ orientation bins $= 128$-d descriptor. High accuracy, robust to illumination and 3D affine viewpoint shifts, but computationally slow.
\textbf{SURF (Bay et al., 2006)}: Replaces Gaussian second derivatives with box filters evaluated via Integral Images; Haar-wavelet responses; 64-d descriptor. $3\times$ faster than SIFT.
\textbf{ORB (Rublee et al., 2011)}: FAST corner detector with intensity centroid orientation (oFAST) $+$ steered BRIEF binary descriptors (rBRIEF). Matches descriptors via bitwise Hamming distance (XOR + popcount) on CPU. Runs $100\times$ faster than SIFT, patent-free, widely used in real-time visual SLAM.}
\Q{State the Optical Flow constraint equation and compare Lucas-Kanade vs Horn-Schunck.}{
\textbf{Brightness Constancy Assumption}: $I(x+u, y+v, t+1) = I(x, y, t)$. First-order Taylor expansion yields:
$$I_x u + I_y v + I_t = 0 \iff \nabla I \cdot \mathbf{v} + I_t = 0$$
This is 1 equation with 2 unknowns $(u, v)$ (Aperture Problem: flow parallel to edges is ambiguous).
\textbf{Lucas-Kanade}: Assumes flow is constant within a small local $n\times n$ window ($n=5$). Solves overdetermined system via least squares:
$\begin{bmatrix} u \\ v \end{bmatrix} = \left(\sum \begin{bmatrix} I_x^2 & I_x I_y \\ I_x I_y & I_y^2 \end{bmatrix}\right)^{-1} \left( -\sum \begin{bmatrix} I_x I_t \\ I_y I_t \end{bmatrix} \right)$.
Fails on large displacements (requires multi-scale image pyramids).
\textbf{Horn-Schunck}: Global variational formulation minimizing brightness error $+$ global smoothness:
$E(u, v) = \iint \left( (I_x u + I_y v + I_t)^2 + \alpha^2 (\|\nabla u\|^2 + \|\nabla v\|^2) \right) dx dy$.
Solves dense flow everywhere via Gauss-Seidel iteration, filling in smooth flow over untextured regions.}
\end{qa}

\begin{trap}[Coordinate \& Geometric Traps]
\begin{itemize}
\item \textbf{Essential vs Fundamental Matrix}: $\mathbf{E}$ operates on normalized ray coordinates $x_{\text{norm}} = \mathbf{K}^{-1} x_{\text{pix}}$ and depends purely on extrinsic pose $[R|t]$. $\mathbf{F}$ operates on raw pixel coordinates and absorbs camera intrinsics $\mathbf{K}$.
\item \textbf{Rank-2 Enforcement in Fundamental Matrix}: Raw linear estimation yields full rank-3 $\mathbf{F}$. If SVD singular values $[\sigma_1, \sigma_2, \sigma_3]$ are not forced to $[\sigma_1, \sigma_2, 0]$, epipolar lines will fail to intersect at a unique epipole.
\end{itemize}
\end{trap}
